%xelatex --shell-escape -interaction=nonstopmode S8-MATE.tex && \
%xelatex --shell-escape -interaction=nonstopmode S8-MATE.tex && \
%xelatex --shell-escape -interaction=nonstopmode S8-MATE.tex

\documentclass[12pt,a4paper,openany,oneside]{book}

% --- Paquetes ---
\usepackage[utf8]{inputenc}
\usepackage[spanish, es-noshorthands]{babel}
\usepackage{amsmath, amsfonts, amssymb}
\usepackage{graphicx}
\usepackage{geometry}
\usepackage{hyperref}
\usepackage{booktabs}
\usepackage{fancyhdr}
\usepackage{listings}
\usepackage{xcolor}
\usepackage{tikz}
\usepackage{pgfplots}
\usepackage{pgfplotstable}
\usepackage{caption}
\usetikzlibrary{babel} % Añade esto a tus librerías de TikZ

\usetikzlibrary{arrows.meta, positioning, shapes.geometric, patterns}
\pgfplotsset{compat=1.18}

\geometry{margin=2.5cm}

% --- Colores y estilos para código Python ---
\definecolor{codegreen}{rgb}{0,0.6,0}
\definecolor{codegray}{rgb}{0.5,0.5,0.5}
\definecolor{codepurple}{rgb}{0.58,0,0.82}
\definecolor{backcolour}{rgb}{0.95,0.95,0.92}

\lstset{
    language=Python,
    backgroundcolor=\color{backcolour},   
    commentstyle=\color{codegreen},
    keywordstyle=\color{blue},
    numberstyle=\tiny\color{codegray},
    stringstyle=\color{codepurple},
    basicstyle=\ttfamily\small,
    breaklines=true,
    frame=single,
    showstringspaces=false
}

% --- Estilo de cabecera y pie ---
\pagestyle{fancy}
\fancyhf{}
\renewcommand{\headrulewidth}{0.4pt}
\renewcommand{\footrulewidth}{0.4pt}
\fancyhead[L]{\small Carlos César Sánchez Coronel}
\fancyhead[R]{\small Sesión 8: Derivadas Parciales y Vector Gradiente}
\fancyfoot[L]{\small Carlos César Sánchez Coronel}
\fancyfoot[C]{\thepage}

% --- Portada ---
\title{
    \textbf{Sesión 8} \\ 
    \Huge \textbf{DERIVADAS PARCIALES Y VECTOR GRADIENTE} \\
    \large La dirección del aprendizaje en espacios multidimensionales \\
    \normalsize Aplicaciones en machine learning, deep learning y visión por computador
}
\author{
    \small Por \\ \vspace{0.2cm}
    \Large \textsc{Carlos César Sánchez Coronel}
}
\date{2026}

\begin{document}

\maketitle
\tableofcontents
\addcontentsline{toc}{chapter}{Índice general}

% ------------------------------------------------------------------
% CAPÍTULO 8: SESIÓN 8 - DERIVADAS PARCIALES Y VECTOR GRADIENTE
% ------------------------------------------------------------------
\chapter{Derivadas Parciales y Vector Gradiente}

En problemas reales de inteligencia artificial, las funciones que manejamos dependen de muchas variables: pesos de una red neuronal, características de entrada, parámetros de un modelo. La derivada ordinaria no es suficiente; necesitamos entender cómo cambia la función cuando modificamos una variable manteniendo las demás constantes. De ahí nacen las \textbf{derivadas parciales}. Al agruparlas obtenemos el \textbf{vector gradiente}, que nos indica la dirección de máximo crecimiento de la función. Este concepto es el corazón del entrenamiento de modelos mediante descenso de gradiente.

\section{Derivadas parciales: sensibilidad individual}

\subsection{Definición intuitiva}
Sea $f(x_1, x_2, \dots, x_n)$ una función de varias variables. La derivada parcial de $f$ respecto a $x_i$ en un punto $\mathbf{a} = (a_1, \dots, a_n)$ mide la tasa de cambio de $f$ cuando solo $x_i$ varía, manteniendo las demás fijas:

\[
\frac{\partial f}{\partial x_i}(\mathbf{a}) = \lim_{h \to 0} \frac{f(a_1, \dots, a_i + h, \dots, a_n) - f(a_1, \dots, a_n)}{h}
\]

Es decir, es la derivada ordinaria de la función $g(t) = f(a_1, \dots, a_{i-1}, t, a_{i+1}, \dots, a_n)$ evaluada en $t = a_i$.

\subsection{Ejemplo básico}
Sea $f(x, y) = x^2 y + \sin(xy)$. Calculamos las derivadas parciales:
\begin{align*}
\frac{\partial f}{\partial x} &= 2xy + y \cos(xy) \\
\frac{\partial f}{\partial y} &= x^2 + x \cos(xy)
\end{align*}

\subsection{Interpretación geométrica}
En el caso de dos variables, $f(x, y)$ representa una superficie en el espacio. La derivada parcial $\frac{\partial f}{\partial x}$ es la pendiente de la curva que se obtiene al cortar la superficie con un plano $y = \text{constante}$. Análogamente para $\frac{\partial f}{\partial y}$.

\begin{center}
\begin{tikzpicture}
\begin{axis}[
    title={Superficie $f(x,y)=x^2 - y^2$},
    xlabel=$x$, ylabel=$y$, zlabel=$f$,
    domain=-2:2,
    y domain=-2:2,
    samples=20,
    mesh/ordering=y varies,
]
\addplot3[surf] {x^2 - y^2};
\addplot3[mesh, domain=-2:2, samples y=1, color=red, thick] (x, 1, {x^2 - 1});
\node[pin=90:{Corte $y=1$}] at (axis cs:1,1,0) {};
\end{axis}
\end{tikzpicture}
\captionof{figure}{La derivada parcial respecto a $x$ es la pendiente de la curva de corte con $y$ constante.}
\end{center}

\subsection{Notación}
Existen varias notaciones para las derivadas parciales:
\begin{itemize}
    \item $\frac{\partial f}{\partial x_i}$ (notación de Leibniz)
    \item $f_{x_i}$ o $f_{i}$ (notación con subíndices)
    \item $D_i f$ (notación de Euler)
\end{itemize}
En el contexto de machine learning, es común encontrar $\frac{\partial \mathcal{L}}{\partial w_{ij}}$ para indicar la derivada de la pérdida respecto a un peso.

\section{El vector gradiente}

\subsection{Definición}
El gradiente de una función $f: \mathbb{R}^n \to \mathbb{R}$ es el vector formado por todas sus derivadas parciales:
\[
\nabla f(\mathbf{x}) = \begin{bmatrix}
\frac{\partial f}{\partial x_1}(\mathbf{x}) \\
\frac{\partial f}{\partial x_2}(\mathbf{x}) \\
\vdots \\
\frac{\partial f}{\partial x_n}(\mathbf{x})
\end{bmatrix}
\]

\subsection{Interpretación geométrica}
\begin{itemize}
    \item El gradiente apunta en la dirección de \textbf{máximo crecimiento} de la función.
    \item Su magnitud $\|\nabla f\|$ indica la tasa de crecimiento en esa dirección.
    \item En dirección perpendicular al gradiente, la función no cambia (primer orden).
\end{itemize}

\subsection{Ejemplo en 2D}
Para $f(x, y) = x^2 + y^2$, el gradiente es $\nabla f = (2x, 2y)$. En el punto $(1,1)$, el gradiente es $(2,2)$, que apunta radialmente hacia afuera (dirección de máximo crecimiento). En la dirección opuesta $(-2,-2)$, la función decrece más rápidamente.

\subsection{Visualización con curvas de nivel}
Las curvas de nivel son líneas donde $f(x,y) = \text{constante}$. El gradiente es perpendicular a estas curvas y apunta hacia donde la función aumenta.

\begin{center}
\begin{tikzpicture}
\begin{axis}[
    title={Curvas de nivel y gradiente de $f(x,y)=x^2+y^2$},
    xlabel=$x$, ylabel=$y$,
    domain=-2:2,
    y domain=-2:2,
    view={0}{90},
    contour gnuplot={levels={0.5,1,2,3,4}, labels=false},
]
\addplot3[contour gnuplot, thick] {x^2 + y^2};
\draw[blue, ->, thick] (axis cs:1,1) -- (axis cs:2,2);
\draw[blue, ->, thick] (axis cs:-1,-1) -- (axis cs:-2,-2);
\end{axis}
\end{tikzpicture}
\captionof{figure}{El gradiente es perpendicular a las curvas de nivel y apunta hacia el crecimiento.}
\end{center}

\subsection{Propiedades importantes}
\begin{enumerate}
    \item \textbf{Linealidad:} $\nabla (\alpha f + \beta g) = \alpha \nabla f + \beta \nabla g$.
    \item \textbf{Regla del producto:} $\nabla (f g) = f \nabla g + g \nabla f$.
    \item \textbf{Regla de la cadena multivariable:} Si $\mathbf{g}(t) = (g_1(t), \dots, g_n(t))$, entonces
    \[
    \frac{d}{dt} f(\mathbf{g}(t)) = \nabla f(\mathbf{g}(t)) \cdot \mathbf{g}'(t)
    \]
\end{enumerate}

\section{Aplicación en IA: Descenso de gradiente}

\subsection{La idea fundamental}
En machine learning, tenemos una función de pérdida $\mathcal{L}(\mathbf{w})$ que depende de los parámetros $\mathbf{w}$ (pesos de la red). Queremos encontrar los pesos que minimicen $\mathcal{L}$. Como el gradiente apunta en la dirección de máximo crecimiento, para minimizar debemos avanzar en la dirección opuesta:
\[
\mathbf{w}_{\text{nuevo}} = \mathbf{w}_{\text{viejo}} - \eta \nabla \mathcal{L}(\mathbf{w}_{\text{viejo}})
\]
donde $\eta$ es la tasa de aprendizaje.

\subsection{Ejemplo paso a paso: regresión lineal con un parámetro}
Consideremos un modelo simplificado $y = wx$ y la pérdida MSE para un solo punto $(x, y)$:
\[
\mathcal{L}(w) = (wx - y)^2
\]
La derivada (gradiente en 1D) es:
\[
\frac{d\mathcal{L}}{dw} = 2(wx - y) \cdot x
\]
El descenso de gradiente actualiza:
\[
w_{\text{nuevo}} = w_{\text{viejo}} - \eta \cdot 2(w_{\text{viejo}} x - y) x
\]

\subsection{Ejemplo con dos parámetros: regresión lineal con intercepto}
Ahora $y = w_1 x + w_2$, pérdida MSE:
\[
\mathcal{L}(w_1, w_2) = (w_1 x + w_2 - y)^2
\]
Las derivadas parciales:
\begin{align*}
\frac{\partial \mathcal{L}}{\partial w_1} &= 2(w_1 x + w_2 - y) \cdot x \\
\frac{\partial \mathcal{L}}{\partial w_2} &= 2(w_1 x + w_2 - y) \cdot 1
\end{align*}
El gradiente es $\nabla \mathcal{L} = \left( \frac{\partial \mathcal{L}}{\partial w_1}, \frac{\partial \mathcal{L}}{\partial w_2} \right)$. La actualización:
\[
w_1 \leftarrow w_1 - \eta \frac{\partial \mathcal{L}}{\partial w_1}, \quad w_2 \leftarrow w_2 - \eta \frac{\partial \mathcal{L}}{\partial w_2}
\]

\subsection{Visualización del proceso}
En la práctica, la función de pérdida puede tener formas complejas. El descenso de gradiente sigue la dirección contraria al gradiente, buscando un mínimo.

\begin{center}
\begin{tikzpicture}
\begin{axis}[
    title={Descenso de gradiente en superficie de pérdida},
    xlabel={$w_1$}, ylabel={$w_2$}, zlabel={$\mathcal{L}$},
    domain=-3:3,
    y domain=-3:3,
    samples=20,
]
\addplot3[surf, opacity=0.5] { (x^2 + y^2)/5 + sin(deg(x))*0.5 };
\addplot3[mesh, domain=-3:3, samples y=1, color=red, thick] 
    coordinates { (2,2,1.5) (1.5,1.5,1.2) (1,1,0.8) (0.5,0.5,0.4) (0,0,0) };
\end{axis}
\end{tikzpicture}
\captionof{figure}{Trayectoria del descenso de gradiente hacia un mínimo.}
\end{center}

\section{Aplicación en Visión por Computador: Detección de bordes}

En una imagen digital, podemos interpretar los niveles de gris como una función $I(x, y)$ (discreta). El gradiente de esta función resalta las regiones donde hay cambios bruscos de intensidad, es decir, los bordes.

\subsection{Gradiente de una imagen}
Para una imagen, el gradiente se aproxima mediante diferencias finitas:
\begin{align*}
\frac{\partial I}{\partial x} &\approx I(x+1, y) - I(x, y) \\
\frac{\partial I}{\partial y} &\approx I(x, y+1) - I(x, y)
\end{align*}
La magnitud del gradiente $G = \sqrt{G_x^2 + G_y^2}$ es alta en los bordes, y la dirección $\theta = \arctan(G_y / G_x)$ indica la orientación del borde.

\subsection{Filtros de Sobel}
Los filtros de Sobel son kernels de convolución que aproximan las derivadas parciales con suavizado:
\[
G_x = \begin{bmatrix}
-1 & 0 & 1 \\
-2 & 0 & 2 \\
-1 & 0 & 1
\end{bmatrix} * I, \quad
G_y = \begin{bmatrix}
-1 & -2 & -1 \\
0 & 0 & 0 \\
1 & 2 & 1
\end{bmatrix} * I
\]



%\begin{center}
%\begin{tikzpicture}
%\draw[step=0.5cm, gray, very thin] (0,0) grid (2,2);
%\foreach \x in {0.25,0.75,1.25,1.75}
%    \foreach \y in {0.25,0.75,1.25,1.75}
%        \fill[gray!\pgfmathparse{(\x+\y)*50}\pgfmathresult] (\x,\y) rectangle ++(0.25,0.25);
%\end{tikzpicture}
%\captionof{figure}{Detección de bordes: el gradiente resalta cambios de intensidad.}
%end{center}

\section{Ejemplos con Python}

\subsection{Cálculo de gradientes con SymPy}
\begin{lstlisting}[language=Python]
import sympy as sp

# Definir variables simbólicas
x, y = sp.symbols('x y')
w1, w2 = sp.symbols('w1 w2')

# Función de pérdida para regresión lineal
loss = (w1*x + w2 - y)**2

# Calcular gradiente
grad_w1 = sp.diff(loss, w1)
grad_w2 = sp.diff(loss, w2)

print("dL/dw1 =", grad_w1)
print("dL/dw2 =", grad_w2)

# Evaluar en un punto
valores = {x: 2, y: 5, w1: 1, w2: 0}
print("dL/dw1 en (2,5,1,0) =", grad_w1.subs(valores))
\end{lstlisting}

\subsection{Visualización de superficie de pérdida y gradiente}
\begin{lstlisting}[language=Python]
import numpy as np
import matplotlib.pyplot as plt
from mpl_toolkits.mplot3d import Axes3D

# Función de pérdida: (w1*x + w2 - y)^2 con x=2, y=5 (fijo)
def loss(w1, w2):
    return (2*w1 + w2 - 5)**2

# Derivadas parciales
def grad_w1(w1, w2):
    return 4*(2*w1 + w2 - 5)

def grad_w2(w1, w2):
    return 2*(2*w1 + w2 - 5)

# Malla para visualización
W1 = np.linspace(-5, 5, 50)
W2 = np.linspace(-5, 5, 50)
W1, W2 = np.meshgrid(W1, W2)
Z = loss(W1, W2)

# Superficie 3D
fig = plt.figure(figsize=(12,5))
ax1 = fig.add_subplot(121, projection='3d')
ax1.plot_surface(W1, W2, Z, cmap='viridis', alpha=0.7)
ax1.set_xlabel('w1')
ax1.set_ylabel('w2')
ax1.set_zlabel('Loss')
ax1.set_title('Superficie de pérdida')

# Curvas de nivel y gradiente en algunos puntos
ax2 = fig.add_subplot(122)
contour = ax2.contour(W1, W2, Z, levels=20, cmap='viridis')
ax2.clabel(contour, inline=True, fontsize=8)

# Dibujar gradiente en algunos puntos
puntos = [(-4,-4), (-2,3), (3,-2), (2,2)]
for w1, w2 in puntos:
    g1 = grad_w1(w1, w2)
    g2 = grad_w2(w1, w2)
    ax2.arrow(w1, w2, g1*0.1, g2*0.1, head_width=0.3, head_length=0.2, fc='red', ec='red')
    ax2.plot(w1, w2, 'ro')

ax2.set_xlabel('w1')
ax2.set_ylabel('w2')
ax2.set_title('Curvas de nivel y vectores gradiente')
plt.tight_layout()
plt.show()
\end{lstlisting}

\subsection{Implementación de descenso de gradiente para regresión lineal}
\begin{lstlisting}[language=Python]
import numpy as np
import matplotlib.pyplot as plt

# Generar datos sintéticos
np.random.seed(42)
X = 2 * np.random.rand(100, 1)
y = 4 + 3 * X + np.random.randn(100, 1)  # y = 4 + 3x + ruido

# Añadir columna de unos para el intercepto
X_b = np.c_[np.ones((100, 1)), X]

# Inicialización
w = np.random.randn(2, 1)
eta = 0.1
n_iter = 1000

# Guardar historia para visualizar
w_history = [w.copy()]

for i in range(n_iter):
    # Calcular gradiente
    grad = 2/len(X) * X_b.T @ (X_b @ w - y)
    # Actualizar pesos
    w = w - eta * grad
    w_history.append(w.copy())

w_optimo = w
print(f"Pesos óptimos: intercepto={w_optimo[0][0]:.4f}, pendiente={w_optimo[1][0]:.4f}")

# Visualizar convergencia
w_history = np.array(w_history).squeeze()
plt.figure()
plt.plot(w_history[:,0], w_history[:,1], 'b.-', alpha=0.5)
plt.plot(w_optimo[0], w_optimo[1], 'ro', markersize=10)
plt.xlabel('Intercepto')
plt.ylabel('Pendiente')
plt.title('Trayectoria del descenso de gradiente')
plt.grid(True)
plt.show()
\end{lstlisting}

\subsection{Detección de bordes con Sobel en una imagen}
\begin{lstlisting}[language=Python]
import numpy as np
import matplotlib.pyplot as plt
from scipy import ndimage
from skimage import data, filters

# Cargar imagen de ejemplo
image = data.camera()  # imagen en escala de grises

# Aplicar filtros Sobel
sx = ndimage.sobel(image, axis=0)  # derivada en x
sy = ndimage.sobel(image, axis=1)  # derivada en y
sobel_magnitud = np.hypot(sx, sy)  # magnitud del gradiente

# Visualizar
fig, axes = plt.subplots(1,3, figsize=(12,4))
axes[0].imshow(image, cmap='gray')
axes[0].set_title('Original')
axes[1].imshow(sobel_magnitud, cmap='gray')
axes[1].set_title('Magnitud del gradiente')
axes[2].imshow(sx, cmap='gray')
axes[2].set_title('Gradiente en x')
plt.tight_layout()
plt.show()
\end{lstlisting}

\section{Notación en papers científicos de IA}

En artículos de machine learning y deep learning, el gradiente aparece en múltiples contextos:

\begin{itemize}
    \item \textbf{Actualización de pesos:} 
    \[
    \theta_{t+1} = \theta_t - \eta \nabla_{\theta} \mathcal{L}(\theta_t)
    \]
    donde $\nabla_{\theta} \mathcal{L}$ denota el gradiente de la pérdida respecto a los parámetros $\theta$.

    \item \textbf{Regla de la cadena en backpropagation:}
    \[
    \frac{\partial \mathcal{L}}{\partial \mathbf{W}^{(l)}} = \frac{\partial \mathcal{L}}{\partial \mathbf{a}^{(l)}} \cdot \frac{\partial \mathbf{a}^{(l)}}{\partial \mathbf{z}^{(l)}} \cdot \frac{\partial \mathbf{z}^{(l)}}{\partial \mathbf{W}^{(l)}}
    \]
    (nótese el uso de derivadas parciales y producto matricial).

    \item \textbf{Gradiente en optimización estocástica:} 
    \[
    \tilde{\nabla} = \frac{1}{|\mathcal{B}|} \sum_{i \in \mathcal{B}} \nabla_{\theta} \ell(\mathbf{x}_i, y_i; \theta)
    \]
    donde $\mathcal{B}$ es un mini-batch.

    \item \textbf{Cálculo de la divergencia en modelos generativos:}
    \[
    \nabla_{\mathbf{x}} \log p(\mathbf{x})
    \]
    (gradiente de la log-verosimilitud, usado en score matching).

    \item \textbf{En visión por computador, la ecuación del flujo óptico:}
    \[
    I_x u + I_y v + I_t = 0
    \]
    donde $I_x$, $I_y$ son las derivadas parciales de la imagen respecto a coordenadas espaciales.
\end{itemize}

También es común encontrar la notación $\nabla$ para el operador nabla, y $\nabla^2$ para el laplaciano (suma de segundas derivadas).

\section{Ejercicios propuestos}

\begin{enumerate}
    \item Calcula las derivadas parciales y el gradiente de $f(x,y,z) = x^2 e^{y} + \ln(xz) - \sin(yz)$.
    \item Para la función de pérdida $\mathcal{L}(\mathbf{w}) = \frac{1}{2} \| \mathbf{X}\mathbf{w} - \mathbf{y} \|_2^2$ (regresión lineal con múltiples variables), demuestra que $\nabla_{\mathbf{w}} \mathcal{L} = \mathbf{X}^\top(\mathbf{X}\mathbf{w} - \mathbf{y})$.
    \item Implementa en Python el descenso de gradiente para la función de Himmelblau: $f(x,y) = (x^2 + y - 11)^2 + (x + y^2 - 7)^2$. Visualiza las curvas de nivel y la trayectoria desde varios puntos iniciales.
    \item Explica por qué el gradiente es perpendicular a las curvas de nivel. Proporciona una demostración intuitiva.
    \item En un paper de redes neuronales, encuentras la expresión $\delta^{(l)} = \nabla_{\mathbf{a}^{(l)}} \mathcal{L} \odot \sigma'(\mathbf{z}^{(l)})$. Identifica cada término y su significado.
\end{enumerate}

\section{Resumen}

En esta sesión hemos aprendido:
\begin{itemize}
    \item Las derivadas parciales como medida de sensibilidad individual de una función multivariable.
    \item El vector gradiente, su interpretación geométrica (dirección de máximo crecimiento) y sus propiedades.
    \item Su aplicación fundamental en machine learning: el descenso de gradiente para minimizar funciones de pérdida.
    \item Su aplicación en visión por computador: detección de bordes mediante gradientes de imagen.
    \item Cómo calcular gradientes simbólicamente con SymPy y visualizarlos en Python.
    \item La notación típica en papers científicos de IA.
\end{itemize}
El gradiente es la herramienta que guía el aprendizaje de casi todos los modelos modernos de inteligencia artificial.

\end{document}